\documentclass[10pt]{article}
\usepackage[utf8]{inputenc}
\usepackage[T1]{fontenc}
\usepackage{amsmath}
\usepackage{amsfonts}
\usepackage{amssymb}
\usepackage[version=4]{mhchem}
\usepackage{stmaryrd}
\usepackage{bbold}
\usepackage{hyperref}
\hypersetup{colorlinks=true, linkcolor=blue, filecolor=magenta, urlcolor=cyan,}
\urlstyle{same}

\title{HOMEWORK 4 MODERN ANALYSIS II - SPRING 2026 }

\author{FRANCESCO LIN}
\date{}


%New command to display footnote whose markers will always be hidden
\let\svthefootnote\thefootnote
\newcommand\blfootnotetext[1]{%
  \let\thefootnote\relax\footnote{#1}%
  \addtocounter{footnote}{-1}%
  \let\thefootnote\svthefootnote%
}

%Overriding the \footnotetext command to hide the marker if its value is `0`
\let\svfootnotetext\footnotetext
\renewcommand\footnotetext[2][?]{%
  \if\relax#1\relax%
    \ifnum\value{footnote}=0\blfootnotetext{#2}\else\svfootnotetext{#2}\fi%
  \else%
    \if?#1\ifnum\value{footnote}=0\blfootnotetext{#2}\else\svfootnotetext{#2}\fi%
    \else\svfootnotetext[#1]{#2}\fi%
  \fi
}

\begin{document}
\maketitle
Problem 1. Prove that

$$
\sum_{n=1}^{\infty} \frac{1}{n^{4}}=\frac{\pi^{4}}{90}
$$

by considering the Fourier series of the $2 \pi$-periodic function $f$ such that $f(x)=|x|$ for $x \in[-\pi, \pi]$. (This is called the triangle wave).

Problem 2. Fix $f \in C^{0}(\mathbb{T})$, and consider the operator

$$
\begin{aligned}
M_{f}:\left(C^{0}(\mathbb{T}),\|\cdot\|_{L^{2}}\right) & \rightarrow\left(C^{0}(\mathbb{T}),\|\cdot\|_{L^{2}}\right) \\
g & \mapsto f g .
\end{aligned}
$$

Prove that this operator is bounded, and has operator norm $\left\|M_{f}\right\|_{\text {op }}=\|f\|_{C^{0}}$.\\
Recall that we showed in Lecture 7 that the Fourier coefficients of a function $f \in C^{0}(\mathbb{T})$ satisfy


\begin{equation*}
|\hat{f}(n)|=o(1 / n) \quad \text { if } f \text { is } C^{1} \tag{1}
\end{equation*}


The following two exercises address other relations between local behavior and decay.\\
Problem 3. Prove that if $f$ is only assumed to be Lipschitz, one can conclude $|\hat{f}(n)|= O(1 / n)$ (i.e. there exists $C>0$ such that $|\hat{f}(n)| \leqslant C / n$ for all $n \neq 0$ ). (For a hint, see Stein-Shakarchi vol I Exercise 15 on p. 91-92).

Problem 4. As a partial converse to (1), prove that any sequence $a_{n}$ such that $\left|a_{n}\right|= O\left(\frac{1}{n^{2+\varepsilon}}\right)$ for some $\varepsilon>0$ is the Fourier series of a $C^{1}$ function. 1 Hint: you might need some results about uniform convergence of derivatives from Modern Analysis I, which you can find for example in Ch. 7 of Baby Rudin.

Department of Mathematics, Columbia University\\
E-mail address: \href{mailto:f12550@columbia.edu}{f12550@columbia.edu}

\footnotetext{${ }^{1}$ Problem 1 shows that you cannot take $\varepsilon=0$; it is a open problem for any $k \in \mathbb{N}$ to give a reasonable characterization of which sequences can arise as Fourier coefficients of $C^{k}$ functions.
}
\end{document}
